
\chapter{Modelagem FrameWeb}
\label{sec-frameweb}
\vspace{-1cm}

\emph{\imprimirtitulo} é um sistema Web cuja arquitetura utiliza \textit{frameworks} comuns no desenvolvimento para esta plataforma. Desta forma, o sistema pode ser modelado utilizando a abordagem FrameWeb~\cite{souza-celebratingfalbo20}.

A Tabela~\ref{tabela-frameworks} indica os \textit{frameworks} presentes na arquitetura do sistema que se encaixam em cada uma das categorias de \textit{frameworks} que FrameWeb dá suporte. Em seguida, os modelos FrameWeb são apresentados para cada camada da arquitetura.

% \vitor{Substituir os valores da segunda coluna da Tabela~\ref{tabela-frameworks} pelos \textit{frameworks} utilizados no seu projeto. Remover o \hl{fundo amarelo}.}


\begin{footnotesize}
	\begin{longtable}{|c|c|}
		\caption{\textit{Frameworks} da arquitetura do sistema separados por categoria.}
		\label{tabela-frameworks}\\\hline
		
		\rowcolor{lightgray}
		\textbf{Categoria de \textit{Framework}} & \textbf{\textit{Framework} Utilizado} \\\hline 
		\endfirsthead
		\hline
		\rowcolor{lightgray}
		\textbf{Categoria de \textit{Framework}} & \textbf{\textit{Framework} Utilizado} \\\hline 
		\endhead

		Controlador Frontal & Lustre \\\hline

		Injeção de Dependências & Wisp \\\hline

		Mapeamento Objeto/Relacional & Squirrel \\\hline

		Segurança & Glow Auth \\\hline
	\end{longtable}
\end{footnotesize}




\section{Camada de Negócio}
\label{sec-frameweb-negocio}

Na camada de negócio, vivem as classes do pacote \textit{Datatypes}, que representam um usuário, que possui nome, email e senha, que pode estar associado a um ou mais feeds RSS, que possuem URL, título, descrição, a data inicial de publicação, a data da última postagem, e um link para uma imagem do feed.

Cada feed pode estar associado a um ou mais itens, que representam cada post do feed, e possuem um título, um link para o post, uma descrição e um link para alguma mídia que esteja ligada ao item, que pode ser, por exemplo, o arquivo de áudio de um episódio de podcast.

Tanto o feed quanto o item podem ter uma ou mais categorias, que possuem o nome dessa categoria, e podem ser extraídas tanto do arquivo do feed RSS quanto ser criadas pelo próprio usuário para que ele possa organizar, filtrar e pesquisar seus feeds.

\section{Camada de Acesso a Dados}
\label{sec-frameweb-dados}

Na camada \textit{Persistence}, vivem as classes do tipo Data Access Object, um para cada tipo da camada de negócio. Nessas classes, além de métodos de criação e apagamento, há também métodos de listagem utilizando critérios específicos, como filtrar por categoria, por feed ou por usuário.

\section{Camada de Apresentação}
\label{sec-frameweb-apresentacao}

Os pacotes restantes lidam com a navegação do \textit{front-end}, ou da comunicação entre o \textit{front-end} com o \textit{back-end}. Existe um pacote para lidar com as páginas de \textit{login} e registro de usuário, e outro para lidar com as páginas para quando o usuário já está logado. Enquanto isso, os pacotes \textit{Control} e \textit{Application} lidam com a comunicação entre as duas pontas, onde o primeiro vive no \textit{front-end} e o segundo \textit{back-end}.